\pagestyle{otFront}
\otPageTitre{Apprendre avec OptiTAB}{\OptionLivre · Terminale générale}
Ce manuel accompagne l'enseignement optionnel français de \OptionLivre.
Il associe un cours progressif, des méthodes expliquées, des exercices de trois niveaux et des solutions rédigées.
\begin{savoirfaire}[Un parcours en quatre temps]
\textbf{Comprendre.} Lire le cours et refaire l'exemple résolu, calculs masqués.\par
\textbf{S'entraîner.} Commencer par les exercices accessibles, puis travailler les applications.\par
\textbf{Chercher.} Réserver un vrai temps aux problèmes, aux conjectures et à leur justification.\par
\textbf{Consolider.} Comparer les raisonnements, noter l'erreur puis refaire la question deux jours plus tard.
\end{savoirfaire}
Les points colorés indiquent une difficulté croissante. Les renvois d'un exercice vers sa solution, et de la solution vers l'énoncé, sont cliquables.
Les six compétences se travaillent ensemble : chercher, modéliser, représenter, raisonner, calculer et communiquer.
\souscours{Lire les énoncés avec précision}
Une hypothèse est une information disponible ; une conclusion est ce qu'il faut établir.
Une représentation ou une simulation aide à conjecturer. Une preuve explique pourquoi l'énoncé reste vrai dans tous les cas annoncés.
Chaque problème exige une conclusion dans son contexte et, si nécessaire, une unité, une précision ou une limite du modèle.
\souscours{Données et outils}
Sauf mention contraire, les données des exercices sont fictives et choisies pour l'apprentissage. Les situations médicales, économiques et cryptographiques sont des modèles pédagogiques.
Une valeur exacte précède autant que possible son arrondi. La calculatrice et Python servent à explorer, à calculer et à contrôler ; les arguments mathématiques restent explicités.
\vfill
{\small OptiTAB · Édition numérique, octobre 2026. Textes, exercices et figures de cette édition : création éditoriale OptiTAB. Les programmes officiels restent la propriété de leurs auteurs institutionnels. Ce manuel indépendant n'est pas une publication du ministère et ne revendique aucune homologation.}
\clearpage
\otPageTitre{Sommaire}{Un itinéraire complet, des outils aux synthèses}
\noindent\makebox[9mm][l]{\otSemi\color{otNavy}01}\hyperref[chap-1]{Suites et modèles discrets\otPoints\pageref*{chap-1}}\par\vspace{2.6mm}
\noindent\makebox[9mm][l]{\otSemi\color{otNavy}02}\hyperref[chap-2]{Exponentielle et logarithme\otPoints\pageref*{chap-2}}\par\vspace{2.6mm}
\noindent\makebox[9mm][l]{\otSemi\color{otNavy}03}\hyperref[chap-3]{Limites, continuité et équations\otPoints\pageref*{chap-3}}\par\vspace{2.6mm}
\noindent\makebox[9mm][l]{\otSemi\color{otNavy}04}\hyperref[chap-4]{Dérivation et optimisation\otPoints\pageref*{chap-4}}\par\vspace{2.6mm}
\noindent\makebox[9mm][l]{\otSemi\color{otNavy}05}\hyperref[chap-5]{Convexité et inégalités\otPoints\pageref*{chap-5}}\par\vspace{2.6mm}
\noindent\makebox[9mm][l]{\otSemi\color{otNavy}06}\hyperref[chap-6]{Primitives et équations différentielles\otPoints\pageref*{chap-6}}\par\vspace{2.6mm}
\noindent\makebox[9mm][l]{\otSemi\color{otNavy}07}\hyperref[chap-7]{Intégrales, aires et valeur moyenne\otPoints\pageref*{chap-7}}\par\vspace{2.6mm}
\noindent\makebox[9mm][l]{\otSemi\color{otNavy}08}\hyperref[chap-8]{Probabilités conditionnelles et Bayes\otPoints\pageref*{chap-8}}\par\vspace{2.6mm}
\noindent\makebox[9mm][l]{\otSemi\color{otNavy}09}\hyperref[chap-9]{Lois discrètes et loi binomiale\otPoints\pageref*{chap-9}}\par\vspace{2.6mm}
\noindent\makebox[9mm][l]{\otSemi\color{otNavy}10}\hyperref[chap-10]{Échantillonnage et décision statistique\otPoints\pageref*{chap-10}}\par\vspace{2.6mm}
\noindent\makebox[9mm][l]{\otSemi\color{otNavy}11}\hyperref[chap-11]{Temps d’attente et lois à densité\otPoints\pageref*{chap-11}}\par\vspace{2.6mm}
\noindent\makebox[9mm][l]{\otSemi\color{otNavy}12}\hyperref[chap-12]{Statistiques à deux variables\otPoints\pageref*{chap-12}}\par\vspace{2.6mm}
\noindent\makebox[9mm][l]{\otSemi\color{otNavy}13}\hyperref[chap-13]{Algorithmique, logique et raisonnement\otPoints\pageref*{chap-13}}\par\vspace{2.6mm}
\vspace{4mm}\repere{dossiers}{Dossiers et évaluations de synthèse}
\repere{corriges-debut}{Corrigés détaillés}
\repere{concordance}{Concordance avec le programme et sources}
\repere{index-methodes}{Index des méthodes et formules}
\clearpage
\otPageTitre{Faire le point}{Diagnostic de départ · À reprendre après quelques semaines}

Sans calculatrice, sauf à la dernière question. Expliquer chaque réponse.
\begin{questions}
\item Simplifier $\frac23+\frac14$, puis $\frac{3^4\times3^2}{3^3}$.
\item Une quantité augmente de $20\,\%$ puis diminue de $20\,\%$. Revient-elle à sa valeur initiale ?
\item Résoudre $x^2-5x+6=0$ et $2x-3\leq5$.
\item Dériver $x^3-2x+1$, $\mathrm e^x$ et $1/x$ sur leurs domaines.
\item Dans un arbre, $P(A)=0{,}3$, $P_A(B)=0{,}8$, $P_{\bar A}(B)=0{,}2$. Calculer $P(B)$.
\item Écrire les quatre premiers termes de $u_0=10$, $u_{n+1}=0{,}8u_n+2$.
\item Calculer à $10^{-3}$ près $\mathrm e^{0{,}4}$ et expliquer ce que produit \code{range(4)} en Python.
\end{questions}
\begin{savoirfaire}[Repères pour se corriger]
1. $11/12$ et $27$. 2. Le coefficient global vaut $1{,}2\times0{,}8=0{,}96$ : baisse de $4\,\%$.
3. $\{2;3\}$ et $x\leq4$. 4. $3x^2-2$, $\mathrm e^x$, $-1/x^2$.
5. $0{,}3\times0{,}8+0{,}7\times0{,}2=0{,}38$.
6. $10,10,10,10$ : la valeur initiale est un équilibre.
7. $1{,}492$ ; \code{range(4)} fournit les entiers $0,1,2,3$.
\end{savoirfaire}
Si les questions 1 à 3 résistent, reprendre les calculs algébriques avant le chapitre 2.
Les questions 4, 5 et 6 orientent respectivement vers les chapitres 4, 8 et 1.

\clearpage\otPageTitre{Mon parcours}{Cocher après une réussite autonome, puis après révision}
\begin{tabularx}{\linewidth}{@{}p{8mm}Yccc@{}}\toprule
& Chapitre & Cours & Exercices & Révision\\\midrule
1 & Suites et modèles discrets & $\square$ & $\square$ & $\square$ \\[2mm]
2 & Exponentielle et logarithme & $\square$ & $\square$ & $\square$ \\[2mm]
3 & Limites, continuité et équations & $\square$ & $\square$ & $\square$ \\[2mm]
4 & Dérivation et optimisation & $\square$ & $\square$ & $\square$ \\[2mm]
5 & Convexité et inégalités & $\square$ & $\square$ & $\square$ \\[2mm]
6 & Primitives et équations différentielles & $\square$ & $\square$ & $\square$ \\[2mm]
7 & Intégrales, aires et valeur moyenne & $\square$ & $\square$ & $\square$ \\[2mm]
8 & Probabilités conditionnelles et Bayes & $\square$ & $\square$ & $\square$ \\[2mm]
9 & Lois discrètes et loi binomiale & $\square$ & $\square$ & $\square$ \\[2mm]
10 & Échantillonnage et décision statistique & $\square$ & $\square$ & $\square$ \\[2mm]
11 & Temps d’attente et lois à densité & $\square$ & $\square$ & $\square$ \\[2mm]
12 & Statistiques à deux variables & $\square$ & $\square$ & $\square$ \\[2mm]
13 & Algorithmique, logique et raisonnement & $\square$ & $\square$ & $\square$ \\[2mm]
\bottomrule\end{tabularx}\par\vspace{5mm}Une case indique une réussite sans aide. Si un exercice résiste, revenir au cours et au premier exemple avant de reprendre le problème.\clearpage\pagestyle{otMain}